\documentclass[11pt, a4paper]{article}

\usepackage[a4paper, margin=2.5cm]{geometry}
\usepackage[utf8]{inputenc}
\usepackage[T1]{fontenc}
\usepackage{helvet} % Opcional, si deseas usar fuente Sans Serif
\renewcommand{\familydefault}{\sfdefault}
\usepackage[spanish, provide=*]{babel}
\babelprovide[import, onchar=ids fonts]{spanish}
%\babelfont{rm}{Noto Sans}
\usepackage{amsmath}
\usepackage{booktabs}
\usepackage{enumitem}

\begin{document}

\begin{center}
    \textbf{\Large TALLER DE REPASO - MATEMÁTICAS 6º}\\[0.2cm]
    \textbf{\large TAU 13: Introducción al Álgebra}
\end{center}
\vspace{0.8cm}

\noindent \textbf{1.
Lee atentamente cada enunciado y escribe la expresión algebraica correspondiente utilizando variables:}
\begin{enumerate}[label=\alph*), itemsep=0.4cm]
    \item La sexta parte de la resta de dos números.
    \item Si a una jarra con un litro de agua se le agregan $x$ litros, ¿cuántos litros de agua tiene en total la jarra?
    \item El triple del cuadrado de un número incrementado en 5.
    \item La razón (división) entre el cubo de un número y el doble de otro.
\end{enumerate}

%\vspace{0.6cm}

\noindent \textbf{2.Para cada una de las siguientes expresiones matemáticas, identifica cuántos términos tiene, cuáles son esos términos, cuáles son sus coeficientes y, finalmente, halla el valor numérico sustituyendo las variables por los números enteros indicados.}
\begin{itemize}
    \item \textbf{Expresión A:} $5x^2 - 3x + 8$ \quad (Evaluar para $x = -2$)
    \item \textbf{Expresión B:} $-4ab^3 + 7a - 1$ \quad (Evaluar para $a = 3, b = -1$)
\end{itemize}

\vspace{0.3cm}
\begin{center}
    \renewcommand{\arraystretch}{1.5}
    \begin{tabular}{l c c c c}
        \toprule
        \textbf{Expresión} & \textbf{Nº de términos} & \textbf{¿Cuáles son?} & \textbf{Coeficientes} & \textbf{Valor Numérico} \\
        \midrule
        \textbf{A} & & & & \\
        \midrule
        \textbf{B} & & & & \\
        \bottomrule
    \end{tabular}
\end{center}

%\vspace{0.6cm}

\noindent \textbf{3. Recuerda que el perímetro de una figura plana es la suma de las longitudes de todos sus lados. Utiliza expresiones algebraicas para resolver:}
\begin{enumerate}[label=\alph*), itemsep=0.4cm]
    \item Un rectángulo tiene un largo representado por la expresión $4x - 1$ y un ancho representado por $2x + 3$. Escribe y simplifica la expresión matemática que representa el perímetro total del rectángulo.
    \item Un triángulo tiene lados que miden $3y$, $5y - 2$ y $y + 4$. Escribe la expresión algebraica simplificada de su perímetro.
    \item Evalúa el perímetro del triángulo del inciso (b) si sabemos que la variable es $y = 3$.
\end{enumerate}

\vspace{0.6cm}

\noindent \textbf{4. Simplifica las siguientes expresiones algebraicas:}
\begin{enumerate}[label=\alph*), itemsep=0.4cm]
    \item $3(4x - 5)2$
    \item $5(2m + 3n - 1)4$
    \item $-2(y^2 - 6y + 3)3$
    \item $4(-x + 7)5$
\end{enumerate}

\end{document}